% MAIN SNAPSHOT — LSPV/RMV::RETURNTRACE1::2026-10-05T05:23:00-07:00
\documentclass[11pt]{article}
\usepackage[margin=.85in]{geometry}
\usepackage{amsmath,amssymb}
\begin{document}
\section*{LSPV-5 RETURNTRACE1: end-to-end SU/vMS return audit}
\textbf{GLOBAL STATUS: OPEN.} P versus NP remains open. No load-bearing implication is closed by this report.

\section*{TARGET}
Attempt the end-to-end trace from an original SAT instance of size T through the vMS parameterization, canonical LDE, RMV line restriction, declared SU05/Hadamard local evaluation, Boolean consumer, and finally an ordinary SAT instance Q. Determine whether a well-typed exponent
\[
|Q_{\rm SAT}|=T^{\zeta+o(1)}
\]
exists in the verified construction.

\section*{CORPUS REVIEWED}
Repository governance, CODEX, pnp-research protocol, root main.tex, TASK-001, MAIN0303R1, RETURNSQUEEZE1, RECURCLOSE1, DELAYOR1, SOURCETRACE1, LOCALBIT1, SEEDBRANCH1, EXTEVAL1, ORBITBATCH1, NORMBLOW1, SHORTKERN1, BATCHSETUP1, CLASSCHARGE1, and LSPV-4 SUCCINCTESCAPE1 were checked. Repository search was run for prior SAT-return, Cook--Levin/Robson, generic-CNF-return, QLEN1, and return-theorem attempts. The causal HEC/CRS-PPC/CCSPC branch and TASK-008 remain parallel and are not strengthened here.

\section*{PRIMARY-LITERATURE PREFLIGHT}
Shaltiel--Umans supplies the algebraic extractor construction. van Melkebeek--Mocelin Sdroievski supplies the targeted-generator/PCPP/reconstruction architecture for AM derandomization. Neither checked source interface supplies a canonical ordinary-CNF SAT-to-smaller-SAT return map from one RMV child. Robson gives a generic RAM-computation-to-SAT encoding of size O(tau log tau); that is an upper-bound encoding, not a contractive source-specific return. Williams' 2024 self-improvement theorem is a positive same-representation precedent for Circuit-SAT and related circuit-analysis problems, not a theorem furnishing the missing ordinary-CNF return here.

\section*{PROVED/EXACT TYPE AUDIT}
For the declared SU05/Hadamard realization, LOCALBIT1 gives
\[
g_i(\widehat g,L,v,u)=\langle x(\widehat g_L(A^iv)),u\rangle_{\mathbb F_2}
\]
and
\[
Q_{L,v,m}=\{L(A^iv):1\le i\le m\}.
\]
ORBITBATCH1 permits source-side reuse of one parent evaluator and removes seed-by-orbit multiplicity as a distinct-address count. NORMBLOW1 gives
\[
N=p^r=T^{\nu+o(1)},\qquad100+o(1)\le\nu\le200+o(1),
\]
so N to the one-half is not a SAT T to the one-half contraction. SUCCINCTESCAPE1 simultaneously forbids charging N merely because the logical object has N coordinates.

The attempted trace reaches a typed interface mismatch after local extraction/reconstruction:
\[
\boxed{\text{the verified source architecture does not presently output an ordinary SAT instance }Q_{\rm SAT}.}
\]
It outputs extractor/reconstruction/PCPP objects and oracle evaluations. No verified source theorem maps one such child to an equisatisfiable ordinary CNF of source-certified sublinear length in T.

Therefore zeta is not merely numerically unknown. In the current architecture it is \textbf{UNDEFINED until a source-matched SAT-return map is specified and proved}. Assigning zeta by identifying the RMV child domain, extractor output, oracle-address set, or Boolean-coordinate list with a SAT instance would be a type error.

\section*{FALSIFICATION OF THE GENERIC REPAIR}
RETURNSQUEEZE1 already proves the generic repair noncontractive. If reconstruction time is
\[
\tau=N^{H_T(K)+o(1)},
\]
a model-matched Cook--Levin/Robson-style return has SAT length tau to the one plus lower-order factors. Calling the same exponent-K SAT decider then costs
\[
N^{K H_T(K)+o(1)}.
\]
For K at least one and H_T(K) at least one, this does not contract K. Generic computation-to-SAT return therefore remains \textbf{FALSIFIED/RETIRED} as the missing Boundary-Squeeze mechanism.

Full-input carry also remains fatal to a direct contraction: if an inherited SAT query explicitly carries T to the one-minus-lower-order bits while the intermediate scale is N=T to the lambda and query length is N to the b, then lambda times b is at least one.

\section*{NEW SHARPENED OPEN ARROW}
The missing object is not compute zeta in isolation. It is a theorem mapping a declared SU/vMS child/reconstruction state to an ordinary SAT instance Q with, in one SAME_REALIZATION: exact satisfiability-preserving semantics; explicit written length; a contractive original-SAT-scale exponent or a typed cycle with product gain below one; charged construction/classification/setup/acquisition; no unresolved current-scale call to the original SAT decider; repeatability; and model matching to the final SAT time-space lower bound.

\section*{HOSTILE-CONTROL AUDIT}
INTERP and RECOMP allow compact evaluators rather than materialization. DELAY forces unresolved current-scale black-box SAT calls to retain the unknown exponent. A maximal semantic oracle can supply source evaluations, so source-side extractor accounting alone relativizes. An algebraic oracle can accelerate LDE/extractor evaluation, so algebraic structure alone does not cross algebrization. Horn, dual-Horn, 2-SAT, affine/Tseitin, bounded-treewidth and independent-product controls prohibit reading compact return structure itself as hardness. COPY/padding cannot manufacture a favorable exponent after normalization.

\section*{DUPLICATION / CONTRADICTION AUDIT}
This run sharpens RETURNSQUEEZE1, NORMBLOW1 and SUCCINCTESCAPE1 into an end-to-end type audit. Earlier wording suggesting that the immediate task was simply to calculate a pre-existing zeta was overstated. The verified corpus contains no source-certified QSAT to measure. SOURCEEDGE1's blanket source-definition absence is stale after SOURCETRACE1/LOCALBIT1, but its no-surrogate rule remains essential.

\section*{WHAT CHANGED RELATIVE TO ROOT MAIN}
The next target should change from compute the SAT-query exponent to construct or falsify the source-specific SAT-return theorem, then compute its exponent. The distinction is load-bearing because the present source architecture terminates at a reconstruction/extractor interface rather than an ordinary SAT instance.

\section*{ADOPTED / REJECTED / RETIRED PRIOR COMPONENTS}
Adopt LOCALBIT1, ORBITBATCH1, NORMBLOW1, SUCCINCTESCAPE1, RETURNSQUEEZE1, DELAYOR1, RECURCLOSE1 and strict SAME_REALIZATION accounting at their declared scopes. Reject identification of RMV child-domain size, extractor-output length, oracle-address count, or Boolean-coordinate count with ordinary SAT-instance length. Retain generic computation-to-SAT return and full-input-carry direct contraction as FALSIFIED/RETIRED mechanisms.

\section*{CAUSAL AND SLICE BRANCHES}
HEC/CRS-PPC/CCSPC remains OPEN at SOURCESTOP1/PATHSTRAT1/OFFPATH1 and still requires an externally fixed non-path-faithful obligation plus an arbitrary-algorithm compiler. TASK-008 remains an independent OPEN monotone-slice route. Neither branch repairs the SAT-return type mismatch proved here.

\section*{SINGLE NEXT TARGET}
State the weakest exact source-specific return relation R_ret(x,S,Q), where S is the declared SU/vMS reconstruction state and Q is an ordinary CNF, and prove
\[
x\in\mathrm{SAT}\iff Q\in\mathrm{SAT}.
\]
Then derive actual written length and construction resources. Falsify first by testing whether every legal Q either carries essentially the full original input, invokes the original SAT decider at current scale, or degenerates to generic computation encoding. Only after such a relation exists is zeta a well-typed quantity.

\section*{FINAL STATUS}
\textbf{PROVED/EXACT:} the present verified SU/vMS source chain does not itself define the ordinary SAT return object whose exponent was requested; zeta is therefore not presently instantiated. \textbf{FALSIFIED/RETIRED:} repairing this gap by generic computation-to-SAT encoding as the contraction mechanism. \textbf{OPEN:} a source-specific satisfiability-preserving return theorem, arbitrary-algorithm acquisition, repeatability, model-matched forbidden-region entry, nonrelativization, nonalgebrization, and P versus NP.

\begin{thebibliography}{5}
\bibitem{SU05} R. Shaltiel and C. Umans, Simple Extractors for All Min-Entropies and a New Pseudorandom Generator, JACM 52(2):172--216, 2005.
\bibitem{vMS25} D. van Melkebeek and N. Mocelin Sdroievski, Instance-Wise Hardness and Refutation versus Derandomization for Arthur--Merlin Protocols, Computational Complexity 34:15, 2025.
\bibitem{Rob91} J. M. Robson, An O(T log T) Reduction from RAM Computations to Satisfiability, Theoretical Computer Science 82(1):141--149, 1991.
\bibitem{Wil24} R. Ryan Williams, Self-Improvement for Circuit-Analysis Problems, STOC 2024.
\bibitem{FLvMV05} L. Fortnow, R. Lipton, D. van Melkebeek and A. Viglas, Time-Space Lower Bounds for Satisfiability, JACM 52(6):835--865, 2005.
\end{thebibliography}
% END_OF_VERIFIED_MAIN::RETURNTRACE1
\end{document}
